\section{The Fear That Saves}

``The fear of the Lord is the beginning of wisdom.''\endnote{Psalm 110:10, Douay--Rheims (Challoner).} The saints who looked at hell did so because they believed that fear, rightly received, is a gift of God; and the Church has confirmed them in it. When it was taught in the sixteenth century that sorrow for sin born of the fear of hell only makes a man a hypocrite, the Council of Trent answered in its decree on penance. It taught that the imperfect contrition called attrition, conceived ``either from the consideration of the turpitude of sin, or from the fear of hell and of punishment,'' if it exclude the will to sin and hope for pardon,

\begin{quote}
not only does not make a man a hypocrite, and a greater sinner, but that it is even a gift of God, and an impulse of the Holy Ghost---who does not indeed as yet dwell in the penitent, but only moves him---whereby the penitent being assisted prepares a way for himself unto justice.\endnote{Council of Trent, session XIV (25 November 1551), \emph{On the Sacrament of Penance}, ch.~4, DS 1678, trans.\ J. Waterworth, \emph{The Canons and Decrees of the Sacred and Oecumenical Council of Trent} (London: Burns and Oates, n.d.; translation first published 1848).}
\end{quote}

\noindent It added an example from Scripture: ``For, smitten profitably with this fear, the Ninevites, at the preaching of Jonas, did fearful penance and obtained mercy from the Lord.'' And among its canons it condemned anyone who should say that the sorrow by which a man thinks over his life, ``pondering on the grievousness, the multitude, the filthiness of his sins, the loss of eternal blessedness, and the eternal damnation which he has incurred, having therewith the purpose of a better life,'' is not a true and profitable sorrow and does not prepare for grace.\endnote{Ibid., canon 5, DS 1705.}

Such fear is not a stain on the Christian life. It is the work of the Holy Spirit in a soul not yet able to love as it should, moving it toward the sacrament in which love is restored. Thomas \`a Kempis says it with homely exactness: ``Nevertheless it is a good thing, if love as yet cannot restrain thee from evil, that at least the fear of hell should hold thee back. But he who putteth aside the fear of God cannot long continue in good, but shall quickly fall into the snares of the devil.''\endnote{\emph{Imitation of Christ} I.24.7.} It is the grace Ignatius teaches the retreatant to ask before the meditation on hell. Aquinas gives the reason God has willed it so. The eternal punishments of the lost, he says, do not heal the lost; but they ``are medicinal punishments for those who refrain from sin through the thought of those punishments.''\endnote{Aquinas, \emph{Summa theologiae} I-II, q.~87, a.~3, reply to objection 2.}

\subsection{The preacher's duty}

For that reason the Church's pastors have never been free to keep silence about hell. Chrysostom, preaching on the fire in which every man's work will be tried, told his people why he would not spare them. He knew, he said, that a chill came over them when they heard of these things; but he could not do otherwise:

\begin{quote}
For this is God's own command, continually to sound these things in your ears\ldots{} and ordained as we have been unto the ministry of the word, we must give pain to our hearers, not willingly but on compulsion.
\end{quote}

\noindent And a little later, with a candour that should comfort every Christian who trembles at the subject:

\begin{quote}
The argument is irksome and pains the hearer: were it only by my own feelings, I know this. For indeed my heart is troubled and throbs; and the more I see the account of hell confirmed, the more do I tremble and shrink through fear. But it is necessary to say these things lest we fall into hell.\endnote{Chrysostom, \emph{Homilies on First Corinthians} 9.1--2, \emph{Nicene and Post-Nicene Fathers}, 1st series, vol.~12 (the Oxford translation, revised by T. W. Chambers).}
\end{quote}

The preacher who speaks of hell is not standing above his hearers. He trembles with them, and he speaks so that none of them, himself included, may fall into it.

\subsection{Neither anxiety nor despair}

The fear the Church commends is not the fear that paralyses. John Paul II, who spoke of hell plainly as ``a real possibility,'' added at once how the thought of it should live in the heart: ``The thought of hell---and even less the improper use of biblical images---must not create anxiety or despair, but is a necessary and healthy reminder of freedom within the proclamation that the risen Jesus has conquered Satan, giving us the Spirit of God who makes us cry `Abba, Father!'\,''\endnote{John Paul II, General Audience of 28 July 1999, n.~4.} A necessary reminder, because freedom is real and can be lost; a healthy one, because it is spoken within the Gospel of the risen Christ, by a Father who wills that we live.

Gregory the Great's deacon, having heard his master's teaching on hell, asked the question every honest reader asks: ``What man is he, though never so holy, that, cometh to leave this mortal life, hath not just cause to fear the unspeakable sentence of damnation?''\endnote{Gregory, \emph{Dialogues} IV.45.} The fear is just. Paul himself commands it: ``with fear and trembling work out your salvation,'' and ``he that thinketh himself to stand, let him take heed lest he fall.''\endnote{Philippians 2:12; 1 Corinthians 10:12, Douay--Rheims (Challoner).} But it is a fear that drives the soul toward God, not away from him. Alphonsus gives the rule for its use in the hour of temptation:

\begin{quote}
When the devil tempts you again to sin, remember Hell, have recourse to God and to the most holy Virgin. The thought of Hell shall preserve you from Hell, because it shall make you have recourse to God.\endnote{Alphonsus, \emph{Preparation for Death}, Consideration XXVIII, third point, p.~347.}
\end{quote}

\noindent That is the whole purpose of the fear: to make us have recourse to God. Its end is love. ``Fear is not in charity: but perfect charity casteth out fear,''\endnote{1 John 4:18, Douay--Rheims (Challoner).} and the \emph{Imitation} adds that ``he who loveth God with all his heart feareth not death, nor punishment, nor judgment, nor hell, because perfect love giveth sure access to God.''\endnote{\emph{Imitation of Christ} I.24.7.} The fear of hell is the servant of that love. It keeps the door while love is still weak, and it gives way, little by little, as love grows strong.
